My earlier work introduced sampling-based planners for \textsl{impulse-based non-prehensile manipulation} (PokeRRT and PokeRRT* \cite{pasricha2022pokerrt}) as a complementary solution to planning with prehensile and quasi-static non-prehensile skills \cite{lynch1999dynamic,zeng2020tossingbot,ruggiero2018nonprehensile,vahrenkamp2010integrated,yu2016more}.
By explicitly modeling impulsive and intermittent contact dynamics through a \textsl{simulation-in-the-loop} approach and constructing a dynamically feasible search tree of poking actions, these planners demonstrated significant advantages over conventional grasping or quasi-static pushing approaches, including vastly expanded reachable workspace, minimal object shape and size restrictions, and inherently faster manipulation.
We further expanded this framework to handle whole-body manipulation tasks under locked, multi-joint failures \cite{lewis1997fault,xie2021maximizing}, leveraging \textsl{precomputed kinodynamic action maps in the robot's task space} that capture the space of impulsive actions available to the constrained system \cite{briscoe2024exploring}.
Experiments show that the use of whole-body impulsive actions significantly expands the robot's failure-constrained reachable workspace while maintaining high manipulation success rates even with unusable end-effectors.
A major limitation of both these works is the computational demand from online, simulation-based re-planning (i.e., exploring impulsive actions for current environment state) that is designed to address the reality gap in frictional contact modeling.